\documentclass[11pt]{article}

\usepackage[margin=1in]{geometry}
\usepackage{array}
\usepackage{booktabs}
\usepackage{longtable}
\usepackage{enumitem}
\usepackage{hyperref}
\usepackage{xcolor}
\usepackage{titlesec}
\usepackage{fancyhdr}
\usepackage{lastpage}
\usepackage{tabularx}
\usepackage{ragged2e}

\hypersetup{colorlinks=true, linkcolor=blue!50!black, urlcolor=blue!50!black}
\setlist[itemize]{topsep=0.25em,itemsep=0.15em,parsep=0em}
\setlist[enumerate]{topsep=0.25em,itemsep=0.15em,parsep=0em}
\titleformat{\section}{\large\bfseries\color{blue!40!black}}{}{0em}{}[\titlerule]
\titleformat{\subsection}{\normalsize\bfseries}{}{0em}{}
\pagestyle{fancy}
\fancyhf{}
\lhead{CPSC 250: Programming for Data Manipulation}
\rhead{Summer 2026}
\cfoot{Page \thepage\ of \pageref{LastPage}}

\newcommand{\coursename}{CPSC 250: Programming for Data Manipulation}
\newcommand{\semester}{Summer 2026}
\newcommand{\instructor}{Dr. Edward J. Brash}
\newcommand{\email}{\href{mailto:edward.brash@cnu.edu}{edward.brash@cnu.edu}}
\newcommand{\office}{Luter 304, but also online!}
\newcommand{\phone}{594-7451}
\newcommand{\officehours}{Online, M--F 1pm--3pm by appointment}
\newcommand{\zybookcode}{CNUCPSC250BrashSummer2026}

\begin{document}

\begin{center}
    {\LARGE \textbf{\coursename}}\\[0.5em]
    {\large \semester}\\[0.75em]
    \emph{This syllabus has been adapted from syllabuses developed previously by Dr. Conner, Dr. Phelps, Professor Perkins, and others.}
\end{center}

\vspace{0.5em}

\begin{tabularx}{\textwidth}{@{}>{\bfseries}l X >{\bfseries}l X@{}}
Course: & CPSC 250 & Instructor: & \instructor \\
Semester: & \semester & Email: & \email \\
Office: & \office & Phone: & \phone \\
Office Hours: & \multicolumn{3}{l@{}}{\officehours} \\
\end{tabularx}

\section{Goals}

Both engineers and scientists need to understand computer programming as well as be able to program computers to solve problems in their fields. This course extends the foundations of computer programming using Python that were introduced in CPSC 150 and focuses on understanding the principles of algorithm development and data structure designs that are central to solving many practical engineering and science problems.

\section{Learning Objectives}

By the end of the course, students should be able to:

\begin{itemize}
    \item Write complete Python programs to solve engineering and science problems using top-down design.
    \item Use high-quality programming standards to develop Python programs to solve practical problems.
    \item Define and use the concepts of data types, arrays, classes, objects, and other complex data structures.
    \item Implement object-oriented programming designs to solve practical problems.
    \item Employ data visualization techniques to display and understand the behavior and underlying nature of complex data sets.
\end{itemize}

\section{Required Text}

We will use Zybooks, a custom text developed by others initially and optimized by me for this course.

\subsection*{Instructions for the class Zybook}

\begin{enumerate}
    \item Sign in or create an account at \href{https://learn.zybooks.com}{learn.zybooks.com}.
    \item Enter zyBook code: \textbf{\zybookcode}.
    \item Subscribe.
\end{enumerate}

A subscription is \$89. I really do apologize for adding extra costs to your education. My only excuse is that I do very much think that this is an exceptional learning resource, and that you will get a lot out of it!

\section{Integrated Development Environment}

My idealistic goal, at this point in time, is to encourage students to work in the operating system with which they are most comfortable: Linux, macOS, or Windows. With that said, we will use \textbf{PyCharm} for the integrated development environment. The underlying Python interpreter will vary depending on the operating system, and I have provided more detailed instructions in supplementary documents.

Students should also expect to use GitHub and git as part of the course workflow. The CPSC 250 course repository will contain lecture notes, examples, and other course materials. The CPSC 250L laboratory repository will be used for lab work and student submissions.

\section{Assessment}

Over my several decades of experience in programming in many different languages, it seems to me that the most important thing that anyone can learn is that:

\begin{quote}
\centering
\textbf{``The only way to learn how to program in a new language is to program in the new language.''}
\end{quote}

With that in mind, you will be asked to continually and continuously, throughout the course, write code in Python. Some of this program writing will be in the form of short code snippets that are designed to reinforce new concepts and ideas. You will also be asked to complete homework and lab assignments that will involve writing complete programs. Finally, there will be tests and a final examination.

In all cases, you will receive credit for this work that will count towards your final grade in the class.

\section{Topics and Daily Content Calendar}

This calendar is extremely tentative and may be adjusted as needed. The course compresses a full semester of material into four weeks, so students should expect a fast pace and should work on programming every day.

\begin{longtable}{@{}>{\raggedright\arraybackslash}p{0.18\textwidth}>{\centering\arraybackslash}p{0.08\textwidth}>{\raggedright\arraybackslash}p{0.66\textwidth}@{}}
\toprule
\textbf{Date} & \textbf{Day} & \textbf{Topic} \\
\midrule
\endfirsthead
\toprule
\textbf{Date} & \textbf{Day} & \textbf{Topic} \\
\midrule
\endhead
Monday 6/1 & 1 & Getting an environment set up; CPSC 150 review \\
Tuesday 6/2 & 2 & Primitive and complex data types and expressions \\
Wednesday 6/3 & 3 & Loops and branching \\
Thursday 6/4 & 4 & Functions, file I/O, and modules \\
Friday 6/5 & 5 & \textbf{Test 1} \\
\addlinespace
Monday 6/8 & 6 & Classes and objects \\
Tuesday 6/9 & 7 & Object-oriented programming I \\
Wednesday 6/10 & 8 & Object-oriented programming II \\
Thursday 6/11 & 9 & Inheritance and polymorphism I \\
Friday 6/12 & 10 & \textbf{Test 2} \\
\addlinespace
Monday 6/15 & 11 & Inheritance and polymorphism II \\
Tuesday 6/16 & 12 & Algorithms I - recursion \\
Wednesday 6/17 & 13 & Algorithms II - searching and sorting \\
Thursday 6/18 & 14 & \textbf{Test 3} \\
Friday 6/19 & -- & \textbf{No class - Juneteenth} \\
\addlinespace
Monday 6/22 & 15 & Data science fundamentals I \\
Tuesday 6/23 & 16 & Data science fundamentals II \\
Wednesday 6/24 & 17 & Data science fundamentals III \\
Thursday 6/25 & 18 & Reading/study day; review for final exam \\
Friday 6/26 & 19 & \textbf{Final Exam} \\
\bottomrule
\end{longtable}

\section{Exams}

There will be three exams and a final exam at the end of the course. The final exam is comprehensive. Test 1 and Test 2 will occur on Friday mornings during class time. Because there are no classes on Friday, June 19, Test 3 is scheduled for Thursday, June 18.

\section{Grading Policy}

Final grades will be based on the following weighting distribution and scale.

\begin{center}
\begin{tabular}{@{}lr@{}}
\toprule
\textbf{Component} & \textbf{Weight} \\
\midrule
Zybooks question sets & 30\% \\
Test 1 & 15\% \\
Test 2 & 15\% \\
Test 3 & 15\% \\
Final Exam & 25\% \\
\bottomrule
\end{tabular}
\end{center}

Final grades will be assigned as follows:

\begin{center}
\begin{tabular}{@{}ll@{}}
\toprule
\textbf{Grade} & \textbf{Percentage Range} \\
\midrule
A & 87--100\% \\
A- & 80--86\% \\
B+ & 77--79\% \\
B & 73--77\% \\
B- & 70--73\% \\
C+ & 67--69\% \\
C & 63--67\% \\
C- & 60--63\% \\
D+ & 57--59\% \\
D & 53--57\% \\
D- & 50--53\% \\
F & below 50\% \\
\bottomrule
\end{tabular}
\end{center}

All Zybooks homework/lab activities are due on the due date. No extensions will be given, except in the case of a valid \textbf{documented} reason.

If you miss a midterm exam because of a valid, \textbf{documented} reason, the grade portion for that exam will be added to your final exam weighting. No make-up exams will be given under any circumstances.

The evaluation of your performance in this course will be based entirely on the regular homework assignments and scheduled exams. There is no possibility to do extra work for extra credit.

\section{Honor Code and Cheating}

As all of you know, CNU has an Honor Code, which we all agreed to by both signing and giving a pledge. The spirit of that Honor Code will be strictly observed in this course.

With that said, we need to have a longer discussion in computer science courses about what ``violating the honor code'' actually looks like. I offer the following, which is adapted from Princeton University's statement on cheating in computer science courses:

Programming is an individual creative process much like composition. You must reach your own understanding of the problem and discover a path to its solution. During this time, discussions with other people are permitted and encouraged. However, when the time comes to write code that solves the problem, such discussions, except with me, are no longer appropriate: the code must be your own work. If you have a question about how to use some feature of Python, or the IDE, or some other relevant application, you can certainly ask your friends or colleagues, but specific questions about code you have written must be treated more carefully.

For each assignment, you must specifically describe, in comments in the code itself, whatever help, if any, that you received from others and tell me the names of any individuals with whom you collaborated. This includes help from friends, classmates, upper-level students, grad students, your Mom, etc.

Do not, under any circumstances, copy another person's code. Incorporating someone else's code into your program in any form is a violation of the Honor Code. This includes adapting solutions or partial solutions to assignments from any offering of this course or any other course. Abetting plagiarism or unauthorized collaboration by ``sharing'' your code is also prohibited. Sharing code in digital form is an especially egregious violation: do not e-mail your code or make your source files available to anyone.

Novices often have the misconception that copying and mechanically transforming a program by rearranging independent code, renaming variables, or similar operations makes it something different. Actually, identifying plagiarized source code is easier than you might think. Not only does plagiarized code quickly identify itself as part of the grading process, but also we can turn to software packages, such as Alex Aiken's renowned MOSS software, for automatic help. Indeed, Zybooks uses MOSS to identify plagiarized code, and gives me amazingly detailed reports. If you cheat, MOSS, and in turn I, will probably catch you!

This policy supplements CNU's academic regulations, making explicit what constitutes a violation for this course.

If you have any questions about these matters, please consult me. Cheating on any work will result in either a score of zero or an F for the course, and/or the filing of a case in the CNU honor court. Violation of the honor code may result in dismissal from the University.

\section{Additional University Information}

The following information is provided by the CNU administration, which you may find useful and with which, in general, I happen to agree. I will have additional comments regarding course materials and intellectual property, and we will discuss these topics in class.

\subsection{University Statement on Diversity and Inclusion}

The Christopher Newport University community engages and respects different viewpoints, understands the cultural and structural context in which those viewpoints emerge, and questions the development of our own perspectives and values, as these are among the fundamental tenets of a liberal arts education.

Accordingly, we affirm our commitment to a campus culture that embraces the full spectrum of human attributes, perspectives, and disciplines, and offers every member of the University the opportunity to become their best self.

Understanding and respecting differences can best develop in a community where members learn, live, work, and serve among individuals with diverse worldviews, identities, and values. We are dedicated to upholding the dignity and worth of all members of this academic community such that all may engage effectively and compassionately in a pluralistic society.

If you have specific questions, suggestions, or concerns regarding diversity on campus please contact \href{mailto:Diversity.Inclusion@CNU.edu}{Diversity.Inclusion@CNU.edu}.

\subsection{Disabilities/Accessibility}

For a student to receive an accommodation due to a disability, that disability must be on record in the Office of Student Affairs, 3rd Floor, David Student Union (DSU). If you have a diagnosed disability, please contact Jacquelyn Barnes, Student Disability Support Specialist in Student Affairs, at 594-7160 to discuss your needs.

Students with documented disabilities are to notify the instructor at least seven days prior to the point at which they require an accommodation, with the first day of class recommended, in private, if accommodation is needed. The instructor will provide students with disabilities with the reasonable accommodations approved and directed by the Office of Student Affairs. Work completed before the student notifies the instructor of his/her disability may be counted toward the final grade at the sole discretion of the instructor.

\subsection{Success}

I want you to succeed in this course and at Christopher Newport University. I encourage you to contact me during office hours or to schedule an appointment to discuss course content or to answer questions you have. If I become concerned about your course performance, attendance, engagement, or well-being, I will contact you first. I also may submit a referral through our Captains Care Program. The referral will be received by the Center for Academic Success as well as other departments when appropriate, including Counseling Services and the Office of Student Engagement. If you are an athlete, the Manager of Athletic Academic Success Programs will be notified. Someone will contact you to help determine what will help you succeed. Please remember that this is a means for me to support you and help foster your success at Christopher Newport University.

\subsection{Public Health}

The university will provide guidance on public health issues and students will be expected to comply with university protocols.

\subsection{Academic Support}

The Center for Academic Success offers free tutoring assistance for Christopher Newport students in several academic areas. Center staff offer individual assistance and/or workshops on various study strategies to help you perform your best in your courses. The center also houses the Alice F. Randall Writing Center. Writing consultants can help you at any stage of the writing process, from invention, to development of ideas, to polishing a final draft. The Center is not a proofreading service, but consultants can help you to recognize and find grammar and punctuation errors in your work as well as provide assistance with global tasks. Contact them as early in the writing process as you can!

You may contact the Center for Academic Success to request a tutor, confer with a writing consultant, obtain a schedule of workshops, or make an appointment to talk with a staff member about study skills and strategies. The Center is located in Christopher Newport Hall, first floor, room 123. You may email \href{mailto:academicsuccess@cnu.edu}{academicsuccess@cnu.edu} or call (757) 594-7684.

\subsection{Course Materials}

All content created and assembled by the faculty member and used in this course is to be considered intellectual property owned by the faculty member and Christopher Newport University. It is provided solely for the private use of the students currently enrolled in this course. To ensure the free and open discussion of ideas, students may not make available any of the original course content, including but not limited to lectures, discussions, videos, handouts, and/or activities, to anyone not currently enrolled in the course without the advance written permission of the instructor. This means that students may not record, download, screenshot, or in any way copy original course material for the purpose of distribution beyond this course. A violation may be considered theft. It is the student's responsibility to protect course material when accessing it outside of the physical classroom space.

\end{document}
